\documentclass[10pt,a4paper]{amsart}
\usepackage[margin=19mm]{geometry}
\usepackage{amsmath,amssymb,mathtools}
\usepackage[scr=rsfs,cal=euler]{mathalfa}
\usepackage{xfrac}
\usepackage[unicode,hidelinks,bookmarks=false]{hyperref}
\newcommand{\ie}{i.e.\ }
\newcommand{\cf}{cf.\ }
\newcommand{\resp}{respectively\ }
\newcommand{\Subsection}{\subsection}
\providecommand{\nobreakdash}{\nobreak\hbox{-}}
\setlength{\emergencystretch}{2em}
\multlinegap0pt
\usepackage{bm}
\usepackage{bbm}
\usepackage{dsfont}
\usepackage{xspace}
\usepackage{cancel}
\usepackage{mathtools}
\usepackage{amsthm}
\makeatletter
\@ifundefined{theorem}{\newtheorem{theorem}{Theorem}}{}
\@ifundefined{lemma}{\newtheorem{lemma}[theorem]{Lemma}}{}
\makeatother
\newcommand{\E}{\mathbb E}
\newcommand{\Pp}{\mathbb P}
\newcommand{\ip}[2]{\left\langle #1,#2\right\rangle}
\newcommand{\norm}[1]{\left\lVert #1\right\rVert}
\newcommand{\one}{\mathbf 1}
\title[Sharp margin-based generalization bounds for realizable SVM]{Sharp margin-based generalization bounds for realizable SVM}
\author{Steve Hanneke, Aryeh Kontorovich}
\date{Source: arXiv:2609.17845v1 (CC BY 4.0)}
\begin{document}
\maketitle

\noindent\emph{Source and licence.} Sharp margin-based generalization bounds for realizable SVM. Version arXiv:2609.17845v1 (CC BY 4.0), distributed under the Creative Commons Attribution 4.0 International licence, as recorded in the arXiv licence metadata for this exact version and shown on the abstract page https://arxiv.org/abs/2609.17845v1. Full rights notice: licenses/arxiv-2609.17845-CC-BY-4.0.txt.\par\smallskip

\noindent\textit{Editorial scope.} Counted unit: the Theorem whose source label is \texttt{thm:entropy} in the article (source0.tex lines 948--966, source label \texttt{thm:entropy}), delivered together with the article's own complete original proof of it (source0.tex lines 968--1122, one \texttt{proof} environment of 155 lines). No mathematics is written, repaired or re-derived: statement and proof are the article's own text line for line. Required context retained uncounted: lem:covariance-identity (lines 666--718); lem:variance-split (lines 786--816); lem:binary-payment (lines 842--864). Every definition, display and label of those lines is retained unchanged. Omitted from this package and not redistributed: 1--665, 719--785, 817--841, 865--947, 967--967, 1123--1934. Nothing inside a retained excerpt is omitted or altered: every retained body is delivered exactly as the article's lines are written, so no omission receipt of any kind accompanies this package. Citations: the retained excerpts carry no citation command at all, so no bibliography entry of the article is transcribed and no reference is redistributed. Reference adaptation: none was needed and none was made; every reference inside the delivered unit resolves inside it, so no unresolved reference or citation is produced. Printed slips kept literal: no printed word, symbol, hypothesis or number of the counted statement or of its proof was altered. Layout and class: the article's own class is \texttt{article} with packages including geometry, amsmath, amssymb, amsthm, mathtools, fontenc, hyperref, microtype, enumitem; the wrapper is the lane's pinned \texttt{amsart} editorial wrapper at 10pt on A4, which restates the theorem-like environments the retained text needs and adds the article's own macro definitions that the retained text needs (\texttt{\textbackslash E}, \texttt{\textbackslash Pp}, \texttt{\textbackslash ip}, \texttt{\textbackslash norm}, \texttt{\textbackslash one}), with extra wrapper packages (bm, bbm, dsfont, xspace, cancel, mathtools). The delivered numbering is the unit's own. Assets: no retained span contains a figure environment, a graphics-inclusion command, a PDF-page inclusion, a table or a TikZ picture; no image file is copied into this package, the package declares no asset dependency and no image file is redistributed with it. Licence: version arXiv:2609.17845v1 of this article is distributed under the Creative Commons Attribution 4.0 International licence, as recorded in the arXiv licence metadata for this exact version and shown on the abstract page https://arxiv.org/abs/2609.17845v1. A full rights notice is delivered at licenses/arxiv-2609.17845-CC-BY-4.0.txt. Characters: every retained excerpt is pure ASCII, so the delivered engine, XeLaTeX with its default Latin Modern text font, reports no missing character.
\begin{lemma}[KKT--covariance identity and an informative coordinate]
\label{lem:covariance-identity}
Let \(\nu\) be admissible at squared-norm budget \(k\). Define
\[
 \bar\alpha_i:=\E\alpha_i(B),
 \qquad
 m_i:=\ip{\bar u}{x_i}.
\]
Then
\[
 \boxed{
 t=
 \sum_{i\in J\cup V}\bar\alpha_i(1-m_i).}
 \tag{3.1}
\]

Assume in addition that every variable coordinate is genuinely random:
\(0<p_i<1\), where
\[
 p_i:=\Pp(i\in B),
 \qquad
 i\in V.
\]
Set
\[
 a_i:=\E[\ip{U}{x_i}\mid i\notin B],
 \qquad
 b_i:=\E[\ip{U}{x_i}\mid i\in B],
 \qquad
 g_i:=a_i-b_i.
\]
Then
\[
 a_i\ge1,
 \qquad
 b_i\le0,
 \qquad
 g_i\ge1,
\]
and
\[
 \boxed{
 t\le
 \sum_{i\in V}\bar\alpha_i p_i g_i.}
 \tag{3.2}
\]
Consequently, if \(t>0\), some \(i\in V\) satisfies
\[
 \boxed{
 p_i g_i\ge\frac{t}{k}.}
 \tag{3.3}
\]
\end{lemma}

\begin{lemma}[Variance drop from revealing one deletion bit]
\label{lem:variance-split}
Fix \(i\in V\) with \(0<p_i<1\), write \(p=p_i\), and put
\[
 \bar u_0:=\E[U\mid i\notin B],
 \qquad
 \bar u_1:=\E[U\mid i\in B],
\]
\[
 t_0:=
 \E[\norm{U-\bar u_0}^2\mid i\notin B],
 \qquad
 t_1:=
 \E[\norm{U-\bar u_1}^2\mid i\in B].
\]
Then
\[
 \boxed{
 t=(1-p)t_0+pt_1+\Delta_i,}
 \qquad
 \Delta_i:=
 p(1-p)\norm{\bar u_0-\bar u_1}^2,
 \tag{3.4}
\]
and
\[
 \boxed{
 \Delta_i\ge p(1-p)g_i^2.}
 \tag{3.5}
\]
\end{lemma}

\begin{lemma}[Binary entropy controlled by the variance drop]
\label{lem:binary-payment}
Let
\[
 0<p\le\frac12,
 \qquad
 0<A\le1,
 \qquad
 g\ge1,
 \qquad
 pg\ge A,
\]
and suppose
\[
 \Delta\ge p(1-p)g^2.
\]
Then
\[
 \boxed{
 h(p)\le4\Delta\log\frac eA.}
 \tag{3.6}
\]
\end{lemma}

\begin{theorem}[Entropy bound for random admissible deletion sets]
\label{thm:entropy}
Let \(\nu\) be any admissible deletion law at squared-norm budget \(k>0\). Then
\[
 \boxed{
 H(B)
 \le
 4\Phi_k(t)+2\E|B|,}
 \qquad
 t=\E\norm{u_B-\E u_B}^2.
 \tag{3.10}
\]
In particular,
\[
 \boxed{
 H(B)\le8k+2\E|B|.}
 \tag{3.11}
\]
\end{theorem}

\begin{proof}
We argue by induction on \(|V|\). The assertion is immediate when \(V\) is
empty.

First remove deterministic coordinates. If \(p_i=0\), then \(i\) is always
retained; move it from \(V\) to \(J\). This changes neither \(B\), \(U\),
\(t\), nor \(\E|B|\). If \(p_i=1\), remove \(i\) from \(V\) and replace
\(B\) by \(B\setminus\{i\}\). Entropy and \(t\) are unchanged, while the
expected cardinality drops by one, so the induction hypothesis for the
reduced law is stronger than the desired conclusion. We may therefore
assume
\[
 0<p_i<1\qquad(i\in V).
 \tag{3.12}
\]

The map \(B\mapsto u_B\) is injective on the support of \(\nu\). Indeed, if
\(B\ne D\), choose \(i\in B\setminus D\), after interchanging the two sets
if necessary. State \(B\) deletes \(i\), whereas state \(D\) retains it, and
hence
\[
 \ip{u_B}{x_i}\le0,
 \qquad
 \ip{u_D}{x_i}\ge1.
\]
Thus \(u_B\ne u_D\). Consequently, if \(\nu\) is not a point mass, then
\(t>0\). The point-mass case has \(H(B)=0\) and is finished, so assume
\(t>0\).

Fix a coordinate \(i\in V\), and reveal the deletion bit
\[
 \xi:=\one_{\{i\in B\}},
 \qquad
 p:=\Pp(\xi=1),
 \qquad
 B':=B\setminus\{i\}.
\]
Conditionally on \(\xi=0\), the coordinate \(i\) becomes permanently retained;
conditionally on \(\xi=1\), it becomes permanently deleted and may be removed
from the ground set. Therefore both conditional laws are admissible at the
same squared-norm budget \(k\) on a variable set of size \(|V|-1\). Let \(t_0,t_1\) be
their separator variances as in Lemma~\ref{lem:variance-split}.

The map \(B\leftrightarrow(\xi,B')\) is one-to-one. Expanding the definition of
Shannon entropy gives the binary chain rule
\[
 H(B)
 =
 h(p)
 +(1-p)H(B'\mid \xi=0)
 +pH(B'\mid \xi=1).
 \tag{3.13}
\]
Also,
\[
 (1-p)\E[|B'|\mid \xi=0]
 +p\E[|B'|\mid \xi=1]
 =
 \E|B|-p.
 \tag{3.14}
\]
Applying the induction hypothesis in the two children and using
(3.13)--(3.14) yields
\[
 \begin{aligned}
 H(B)
 &\le
 h(p)
 +4\left(
      (1-p)\Phi_k(t_0)
      +p\Phi_k(t_1)
   \right)
 +2(\E|B|-p).
 \end{aligned}
 \tag{3.15}
\]
We now choose the splitting coordinate.

Suppose first that some coordinate has \(p>1/2\), and split on any such
coordinate. Since \(\Phi_k\) is concave and increasing, and since
Lemma~\ref{lem:variance-split} gives
\[
 (1-p)t_0+pt_1=t-\Delta_i\le t,
\]
we have
\[
 (1-p)\Phi_k(t_0)+p\Phi_k(t_1)
 \le
 \Phi_k(t-\Delta_i)
 \le
 \Phi_k(t).
 \tag{3.16}
\]
Moreover,
\[
 h(p)\le\log2<1<2p.
 \tag{3.17}
\]
Substitution into (3.15) proves (3.10).

It remains to consider the case
\[
 p_i\le\frac12\qquad(i\in V).
 \tag{3.18}
\]
By Lemma~\ref{lem:covariance-identity}, choose \(i\) so that
\[
 p_i g_i\ge\frac tk.
\]
Set
\[
 A:=\frac tk\in(0,1].
\]
Lemmas~\ref{lem:variance-split} and~\ref{lem:binary-payment} give
\[
 h(p_i)
 \le
 4\Delta_i\log\frac{ek}{t}.
 \tag{3.19}
\]
Concavity of \(\Phi_k\) and (3.4) imply
\[
 (1-p_i)\Phi_k(t_0)+p_i\Phi_k(t_1)
 \le
 \Phi_k(t-\Delta_i).
 \tag{3.20}
\]
Because a differentiable concave function lies below each of its tangent
lines,
\[
 \Phi_k(t)-\Phi_k(t-\Delta_i)
 \ge
 \Delta_i\Phi_k'(t)
 =
 \Delta_i\log\frac{ek}{t}.
 \tag{3.21}
\]
Combining (3.19)--(3.21),
\[
 h(p_i)
 +4\left(
      (1-p_i)\Phi_k(t_0)
      +p_i\Phi_k(t_1)
   \right)
 \le
 4\Phi_k(t).
\]
Equation (3.15) now proves (3.10), completing the induction.

Finally,
\[
 t\le\E\norm U^2\le k
\]
and \(\Phi_k(t)\le2k\), so (3.11) follows.
\end{proof}

\end{document}
